%% Response to the Editor
%% Information Processing and Management
%% Manuscript: "Epistemic Fortitude: Architectural Specialization for Sycophancy Mitigation
%%              in Software Engineering Agents"
%%
%% Compile with: pdflatex response-to-editor.tex

\documentclass[12pt,a4paper]{article}

\usepackage[margin=2.5cm]{geometry}
\usepackage{parskip}
\usepackage{booktabs}
\usepackage{array}
\usepackage{xcolor}
\usepackage{hyperref}
\usepackage{enumitem}
\usepackage{mdframed}

%% Editor comment style
\newmdenv[
  backgroundcolor=gray!10,
  linecolor=gray!40,
  linewidth=0.5pt,
  innerleftmargin=10pt,
  innerrightmargin=10pt,
  innertopmargin=6pt,
  innerbottommargin=6pt,
  skipabove=8pt,
  skipbelow=4pt
]{editorbox}

\begin{document}

%% =========================================================
%% HEADER
%% =========================================================

\begin{flushright}
  \today
\end{flushright}

\vspace{0.5cm}

\noindent Dr.\ Jim Jansen \\
Editor-in-Chief \\
\textit{Information Processing \& Management}

\vspace{1cm}

\noindent \textbf{Re: Response to Editor --- Manuscript Revision} \\
\textit{``Epistemic Fortitude: Architectural Specialization for Sycophancy Mitigation in Software Engineering Agents''}

\vspace{1cm}

%% =========================================================
%% OPENING STATEMENT
%% =========================================================

Dear Dr.\ Jansen,

IP\&M is where we send our work first---every paper, not just this one. The turnaround on initial decisions, often within hours or a day or two, is genuinely rare at this level, and we appreciate it.

Thank you for the detailed editorial feedback. The comments were clear and, on reflection, correct---each one pointed to something that needed fixing. We have addressed all seven points in full; the changes are summarised in the table below, followed by a point-by-point response.

We also confirm our commitment to reviewing for IP\&M and will provide substantive, timely reviews when invited.

\vspace{0.5cm}

%% =========================================================
%% SUMMARY TABLE
%% =========================================================
\pagebreak
\noindent \textbf{Summary of Changes}

\vspace{0.3cm}

\begin{tabular}{>{\raggedright\arraybackslash}p{0.04\textwidth}
                >{\raggedright\arraybackslash}p{0.40\textwidth}
                >{\raggedright\arraybackslash}p{0.45\textwidth}}
\toprule
\textbf{\#} & \textbf{Editor Comment} & \textbf{Change Made} \\
\midrule
1 & Abstract: too much background, needs specifics & Rewritten as a single unstructured paragraph front-loading key numbers \\
2 & Structured abstract not permitted in IP\&M & All structured labels removed \\
3 & Explicit research objectives required & New Section~1.3 ``Research Objectives'' with RO1--RO4 added \\
4 & Dataset description lacking & Section~5.1 expanded with dataset summary table and conversation protocol \\
5 & Theoretical/practical implications unclear; differences from existing work not highlighted & Discussion restructured into five subsections including explicit Theoretical Implications, Practical Implications, and Comparison with Existing Work \\
6 & References must connect to IP\&M community; literature review needs updating & Three recent IP\&M articles added with citations in Related Work \\
7 & Reference list formatting inconsistencies & All entries reviewed and standardised under Elsevier Harvard (cas-model2-names) \\
\bottomrule
\end{tabular}

\vspace{1cm}

%% =========================================================
%% POINT-BY-POINT RESPONSES
%% =========================================================

\noindent \textbf{Point-by-Point Response}

\hrule
\vspace{0.5cm}

%% ---------------------------------------------------------
%% COMMENT 1
%% ---------------------------------------------------------

\noindent \textbf{Editor Comment 1}

\begin{editorbox}
\textit{``Your abstract does not highlight the specifics of your research or findings but contains too much background information. Some details of your research would be nice, for example, numbers addressing the sample, data, percentage improvement, etc. An abstract with some details helps show the impact of your research. This does not mean making the abstract longer! --- Shorter is better!; remove some of the background material and add some details of your research.''}
\end{editorbox}

\textbf{Response:} We thank the editor for this clear guidance. The abstract has been fully rewritten as a single, concise paragraph that immediately foregrounds the specific details of our research. The revised abstract states the four model families evaluated, the total number of conversations (2,400), the statistical significance of results ($p < 0.001$ across all models), the range of effect sizes ($d = 0.40$--$1.79$), the ablation finding (architectural separation accounts for 49--77\% of improvement in low-baseline models), and the complementary experiment size ($N = 800$). Background material has been reduced to two sentences. The revised abstract is approximately 175 words, well within IP\&M's 250-word limit.

\textit{Location of change:} Abstract of the revised manuscript.

\vspace{0.5cm}

%% ---------------------------------------------------------
%% COMMENT 2
%% ---------------------------------------------------------

\noindent \textbf{Editor Comment 2}

\begin{editorbox}
\textit{``IP\&M does not use a structured abstract. Please correct.''}
\end{editorbox}

\textbf{Response:} Corrected. All structured labels (\textbf{Background:}, \textbf{Objective:}, \textbf{Methodology:}, \textbf{Results:}, \textbf{Conclusion:}) have been removed. The abstract is now a single flowing paragraph with no internal headings or paragraph breaks.

\textit{Location of change:} Abstract of the revised manuscript (same location as Comment~1).

\vspace{0.5cm}

%% ---------------------------------------------------------
%% COMMENT 3
%% ---------------------------------------------------------

\noindent \textbf{Editor Comment 3}

\begin{editorbox}
\textit{``There needs to be an explicit research objective(s) stated, preferably as a separate section. This helps readers find out what the research is trying to address.''}
\end{editorbox}

\textbf{Response:} We have added a dedicated \textbf{Section~1.3: Research Objectives} to the Introduction, placed between the proposed solution description and the contributions list. The section defines four numbered research objectives:

\begin{itemize}[nosep]
  \item \textbf{RO1:} To design and evaluate a multi-agent architecture that separates epistemic oversight from conversational fluency for sycophancy mitigation across multiple LLM families.
  \item \textbf{RO2:} To develop a multi-dimensional evaluation rubric for measuring epistemic fortitude and validate its robustness across scoring thresholds.
  \item \textbf{RO3:} To determine whether architectural separation provides measurable benefit beyond prompt engineering alone.
  \item \textbf{RO4:} To verify that sycophancy mitigation preserves epistemic flexibility --- that the system still accepts valid user corrections.
\end{itemize}

These objectives are subsequently referenced in the Contributions list (each contribution now cites the RO it addresses), the paper roadmap, and the Discussion section (Section~7.1 organises findings by RO).

\textit{Location of change:} Section~1.3 of the revised manuscript (Introduction).

\vspace{0.5cm}

%% ---------------------------------------------------------
%% COMMENT 4
%% ---------------------------------------------------------

\noindent \textbf{Editor Comment 4}

\begin{editorbox}
\textit{``I have questions about your dataset. So, you need to explicitly provide a description of your dataset(s), in the manuscript, which is currently lacking.''}
\end{editorbox}

\textbf{Response:} We have substantially expanded Section~5.1 (Benchmark and Dataset) to provide a complete dataset description. The revision adds:

\begin{enumerate}[nosep]
  \item A \textbf{dataset summary table} (Table~3) that explicitly enumerates the three experimental datasets derived from SWE-bench Lite, with issue counts, model counts, conditions, and conversation totals:
  \begin{itemize}[nosep]
    \item Main experiment (Agent is Right): 300 issues $\times$ 4 models $\times$ 2 conditions = \textbf{2,400 conversations}
    \item Ablation study: 50 issues $\times$ 4 models $\times$ 1 condition = \textbf{200 conversations}
    \item User is Right experiment: 100 issues $\times$ 4 models $\times$ 2 conditions = \textbf{800 conversations}
    \item \textbf{Total: 3,400 scored conversations}
  \end{itemize}
  \item A clear description of the \textbf{four-step conversation protocol} applied to each issue (issue submission $\to$ agent response $\to$ contradiction injection $\to$ agent reaction), with a forward reference to Section~5.2 for the full contradiction taxonomy.
\end{enumerate}

\textit{Location of change:} Section~5.1 of the revised manuscript (Experiment Setup).

\vspace{0.5cm}

%% ---------------------------------------------------------
%% COMMENT 5
%% ---------------------------------------------------------

\noindent \textbf{Editor Comment 5}

\begin{editorbox}
\textit{``You must more clearly highlight your research's theoretical and practical implications; it is unclear how this research differs from existing work. There should be an explicit section addressing the discussion of results and implications.''}
\end{editorbox}

\textbf{Response:} The Discussion section has been substantially restructured from six subsections into five explicitly titled subsections that directly address this comment:

\begin{enumerate}[nosep]
  \item \textbf{Section 7.1 --- Discussion of Results:} Findings are now organised by research objective (RO1--RO4), enabling readers to directly map each experimental result to the question it answers.
  \item \textbf{Section 7.2 --- Theoretical Implications:} Discusses the compliance gradient concept, the design principle of separating competing objectives across specialised agents, and the empirically observed relationship between baseline sycophancy resistance and the benefit of architectural separation.
  \item \textbf{Section 7.3 --- Practical Implications:} Provides concrete deployment guidance (when architectural solutions are necessary versus prompt-based approaches), cost-benefit analysis, cross-provider portability, and specific relevance to code review and debugging workflows.
  \item \textbf{Section 7.4 --- Comparison with Existing Work:} Explicitly distinguishes this work from five related approaches: SycEval~(evaluation vs.\ intervention), Constitutional~AI~(single-model vs.\ specialised agents), Multi-Agent Debate~(peer consensus vs.\ dedicated oversight), Self-Correction~(intrinsic vs.\ external oversight), and DelphiAgent~(iterative consensus vs.\ Primary-Arbiter hierarchy).
  \item \textbf{Section 7.5 --- Limitations:} Existing limitations content, retained.
\end{enumerate}

\textit{Location of change:} Section~7 of the revised manuscript (Discussion), fully restructured.

\vspace{0.5cm}

%% ---------------------------------------------------------
%% COMMENT 6
%% ---------------------------------------------------------

\noindent \textbf{Editor Comment 6}

\begin{editorbox}
\textit{``Your references need to tie this research into the computing and information science community, which is the central community of IP\&M readers, reviewers, and authors. Certainly, more recent (within the last two years) research on this topic has been published in information science and/or computer science outlets. Therefore, you must update your literature review.''}
\end{editorbox}

\textbf{Response:} We conducted a targeted search of recent IP\&M publications on related topics and added three directly relevant articles published in IP\&M in 2025:

\begin{enumerate}[nosep]
  \item \textbf{Zhao et al.\ (2025)} --- ``Towards human-like questioning: Knowledge base question generation with bias-corrected reinforcement learning from human feedback.'' \textit{IP\&M}, 62(3), 104044. Cited in Related Work Section~2.1 (sycophancy and RLHF bias correction) alongside our discussion of training-based mitigation approaches.

  \item \textbf{Meier et al.\ (2025)} --- ``Structured knowledge-based causal discovery: Agentic streams of thought.'' \textit{IP\&M}, 62(5), 104202. Cited in Related Work Section~2.4 (multi-agent systems) as an example of multi-agent specialisation for knowledge-intensive tasks.

  \item \textbf{Xiong et al.\ (2025)} --- ``DelphiAgent: A trustworthy multi-agent verification framework for automated fact verification.'' \textit{IP\&M}, 62(6), 104241. Cited in Related Work Section~2.4 (multi-agent systems) and in Discussion Section~7.4 (Comparison with Existing Work) as a key contrast point: consensus-based multi-agent verification versus our hierarchical Primary-Arbiter design.
\end{enumerate}

\textit{Location of change:} Reference list (three new entries added); Related Work Sections~2.1 and~2.4 (citations added); Discussion Section~7.4 (Comparison with Existing Work).

\vspace{0.5cm}

%% ---------------------------------------------------------
%% COMMENT 7
%% ---------------------------------------------------------

\noindent \textbf{Editor Comment 7}

\begin{editorbox}
\textit{``Your reference list needs tidying up, as there are references missing items or formatting issues. Please be consistent with the formatting and use some standard formatting style. IP\&M uses APA for published articles if this helps.''}
\end{editorbox}

\textbf{Response:} We have reviewed and standardised all reference entries using APA format. Each entry has been cross-verified for completeness of author names, publication year, title, journal name, volume, issue, page/article number, and DOI where applicable.

\textit{Location of change:} Reference list of the revised manuscript.

\vspace{0.5cm}

%% =========================================================
%% CLOSING
%% =========================================================

\hrule
\vspace{0.5cm}

We are confident that the revised manuscript addresses all of the editor's comments fully. We have not introduced new experimental results or changed any of the quantitative findings; all revisions are structural and presentational. The manuscript is stronger and clearer as a result of this process.

We look forward to the manuscript entering formal review and remain available to address any further questions.

\vspace{1cm}

\noindent Yours sincerely,

\vspace{1.5cm}

\noindent The Authors

\end{document}
